\documentclass[11pt]{article}

\usepackage[utf8]{inputenc}
\usepackage[T1]{fontenc}
\usepackage[margin=1in]{geometry}
\usepackage{booktabs}
\usepackage{array}
\usepackage{xcolor}
\usepackage{fancyvrb}
\usepackage[hidelinks]{hyperref}
\usepackage{url}

\newcolumntype{L}[1]{>{\raggedright\arraybackslash}p{#1}}

\title{GPU-Aware Autoscaling in Cloud Native AI Infrastructure:\\
A Node-Level Telemetry Approach with KEDA and NVML}

\author{Pavan Madduri\\
Independent Researcher\\
\texttt{https://orcid.org/0009-0007-3795-7593}}

\date{September 2026}

\begin{document}
\maketitle

\begin{abstract}
Kubernetes treats GPUs as opaque, countable resources: a pod is either allocated
a device or it is not, with no notion of utilization, memory pressure, or thermal
state. This is acceptable for CPU-bound services but a poor fit for AI inference,
where a pod can report single-digit CPU usage while its GPU is fully saturated.
The Horizontal Pod Autoscaler, which reacts to CPU and memory, therefore cannot
scale GPU workloads on the signal that actually matters. The common workaround
chains an exporter, Prometheus, a query language, and an adapter, which adds
operational complexity and tens of seconds of metric latency. We describe
\texttt{keda-gpu-scaler}, an open-source KEDA external scaler that reads GPU
metrics directly from the NVIDIA Management Library (NVML) on each node and
serves them to Kubernetes autoscaling over gRPC, removing the intermediate
metrics pipeline. We explain why GPU support must run as a DaemonSet rather than
inside the KEDA operator, present the architecture and scaling profiles for
inference engines such as vLLM and Triton, and report production observations
from a multi-node A100 cluster in which direct NVML reads reduced metric latency
from the 15--30~s typical of the Prometheus pipeline to 2--4~s while supporting
GPU scale-to-zero.
\end{abstract}

\section{Introduction}

Running GPU workloads on Kubernetes exposes a gap in the platform's autoscaling
story. The Horizontal Pod Autoscaler (HPA) observes CPU and memory. A large
language model (LLM) serving pod can be handling hundreds of concurrent requests
with its GPU fully saturated while reporting under 10\% CPU, and the HPA takes no
action. The Vertical Pod Autoscaler does not address this either. KEDA
\cite{keda} brings custom-metric scaling within reach, but the standard route to
GPU metrics---an exporter feeding Prometheus, a PromQL query, and either a
Prometheus adapter or KEDA's Prometheus scaler---adds latency and several
components to operate.

\texttt{keda-gpu-scaler} \cite{kgs} takes a different approach: it reads GPU
metrics directly from NVML \cite{nvml} on each node and serves them to KEDA over
gRPC, skipping the metrics pipeline. It runs as a DaemonSet that polls NVML on a
short interval and implements KEDA's external scaler interface
\cite{keda-extscaler}. It ships scaling profiles for inference and training
workloads and handles multi-GPU nodes with configurable aggregation. This paper
describes the problem, why native integration into KEDA is not viable, the
architecture, and an evaluation based on production use.

\section{The GPU Scaling Blind Spot}

\subsection{Current state}
Kubernetes models GPUs as extended resources (e.g., \texttt{nvidia.com/gpu}). A
workload requests a count and the scheduler assigns devices. There is no
supported way for an autoscaler to ask how busy a GPU is or how much device
memory remains. To the HPA, an idle GPU and a saturated GPU are indistinguishable.

\subsection{The metrics pipeline problem}
The conventional path to GPU-based scaling chains multiple components, for
example \texttt{DCGM Exporter $\rightarrow$ Prometheus $\rightarrow$ PromQL
$\rightarrow$ Prometheus Adapter $\rightarrow$ HPA}, or a KEDA variant ending in
KEDA's Prometheus scaler. This works but has drawbacks, summarized in
Table~\ref{tab:pipeline}.

\begin{table}[h]
\centering
\begin{tabular}{L{0.32\textwidth} L{0.58\textwidth}}
\toprule
\textbf{Problem} & \textbf{Impact} \\
\midrule
Metric latency & 15--30~s from GPU state change to scaling decision, from scrape intervals and aggregation \\
Operational complexity & Five components to deploy, configure, and maintain \\
PromQL fragility & Scaling logic encoded in query strings; no type safety, hard to test \\
Single point of failure & A Prometheus outage stops GPU scaling \\
Resource overhead & Prometheus retention, a per-node exporter, and an adapter deployment \\
\bottomrule
\end{tabular}
\caption{Drawbacks of the conventional metrics-pipeline approach.}
\label{tab:pipeline}
\end{table}

\subsection{Why this matters now}
GPU workloads do not scale like CPU workloads. LLM inference has strict latency
targets, so scaling late drops requests. Device memory, rather than compute, is
often the bottleneck: engines such as vLLM \cite{vllm} pre-allocate a key-value
(KV) cache, so memory pressure is the leading signal of queue growth. GPUs are
also expensive, so over-provisioning a handful of accelerators can cost more than
over-provisioning many CPU cores. Finally, idle development and batch inference
endpoints should release GPUs entirely, which standard HPA cannot do for
device-backed pods.

\section{Why Native Integration Fails}

Embedding GPU support directly in the KEDA operator is not viable, for three
reasons.

\paragraph{CGO constraint.} NVIDIA's Go bindings \cite{gonvml} call into
\texttt{libnvidia-ml.so} through cgo. KEDA builds its operator with
\texttt{CGO\_ENABLED=0} to ship static binaries; introducing a cgo dependency
would break that build and release model.

\paragraph{Node-level hardware access.} NVML reads GPU state through device files
(\texttt{/dev/nvidiactl}, \texttt{/dev/nvidia0..N}) that exist only on the
physical GPU node. The KEDA operator runs as a single centralized Deployment and
has no access to devices on worker nodes. A DaemonSet is the appropriate pattern:
each instance runs on a GPU node, mounts the device files, and reads metrics
locally.

\paragraph{Independent release cycle.} GPU tooling evolves quickly---new
architectures, new inference servers, and new metrics such as Multi-Instance GPU
(MIG) partitions. KEDA releases coordinate across many scalers, so shipping GPU
changes on that cadence would be slow. A standalone component can iterate
independently. This design was discussed in the KEDA project \cite{keda7538}.

\section{Architecture}

\subsection{Overview}
The scaler runs as a DaemonSet on each GPU node and serves KEDA over gRPC:

\begin{Verbatim}[fontsize=\small]
GPU Node                                    KEDA Operator
+------------------------------+           +------------------+
|  DaemonSet: keda-gpu-scaler  |           |                  |
|                              |           |  ExternalScaler  |
|  +------------+              |  gRPC     |  trigger config  |
|  | NVML poller|--metrics-->  |--:6000--> |                  |
|  | (2s loop)  |              |           |  -> HPA decision  |
|  +------------+              |           |  -> scale up/down |
|       ^                      |           +------------------+
|  libnvidia-ml.so             |
|  /dev/nvidia0..N             |
+------------------------------+
\end{Verbatim}

\subsection{Data flow}
Each poll cycle (2~s by default) reads streaming-multiprocessor (SM) utilization,
memory-controller utilization, device memory used and total, temperature, and
power draw. Metrics are cached in memory with no disk or external store. When KEDA
calls \texttt{GetMetrics} over gRPC, the scaler returns the requested metric using
the aggregation method named in the \texttt{ScaledObject}, and KEDA feeds the
value to the HPA for a scale up, down, or to-zero decision.

\subsection{gRPC interface}
The scaler implements four methods of KEDA's external scaler contract:
\texttt{IsActive} (true when a metric exceeds the activation threshold, enabling
scale-from-zero), \texttt{StreamIsActive} (a streaming form for push-based
activation), \texttt{GetMetricSpec} (metric name and target for the HPA), and
\texttt{GetMetrics} (the current value). Using gRPC and protobuf gives type
safety, avoids PromQL string parsing, and keeps Prometheus out of the critical
path.

\subsection{Security model}
The DaemonSet needs only read-only access to NVIDIA device files and no
cluster-wide RBAC. The gRPC port is exposed as a ClusterIP Service and is not
reachable outside the cluster. The optional Prometheus endpoint can be disabled.
No secrets are required, and all NVML calls are read-only.

\section{Scaling Profiles for AI Workloads}

Rather than requiring operators to reason about what ``80\% GPU utilization''
means for a given workload, the scaler provides profiles with sensible defaults
(Table~\ref{tab:profiles}); any field can be overridden in the
\texttt{ScaledObject}.

\begin{table}[h]
\centering
\begin{tabular}{L{0.20\textwidth} L{0.20\textwidth} l l L{0.26\textwidth}}
\toprule
\textbf{Profile} & \textbf{Primary metric} & \textbf{Thr.} & \textbf{Act.} & \textbf{Rationale} \\
\midrule
\texttt{vllm-inference} & VRAM utilization & 80\% & 5\% & KV cache pressure tracks queue growth \\
\texttt{triton-inference} & SM utilization & 75\% & 10\% & SM saturation degrades throughput \\
\texttt{training} & SM utilization & 90\% & 0\% & Avoid killing checkpoints; no scale-to-zero \\
\texttt{batch} & VRAM utilization & 70\% & 1\% & Aggressive scale-down to release GPUs \\
\bottomrule
\end{tabular}
\caption{Default scaling profiles.}
\label{tab:profiles}
\end{table}

\section{Multi-GPU Aggregation}

On nodes with four or eight GPUs, per-device metrics must be reduced to a single
value for KEDA. The scaler supports \texttt{max} (default; scale when any GPU
hits threshold, best for inference where one hot device signals a latency
problem), \texttt{avg} (best for training where devices should load evenly),
\texttt{min} (conservative), \texttt{sum} (capacity-based), and the \texttt{p95}
and \texttt{p99} percentiles.

\section{NUMA-Aware GPU Scheduling}

The scaler decides \emph{when} to add pods; where those pods land is a scheduling
concern that affects performance. In a dual-socket server each GPU is wired to one
NUMA node over PCIe. If a pod's CPU is placed on one NUMA node and its GPU on
another, memory accesses cross the inter-socket link, adding latency that
accumulates during model loading and inference. Topology-aware placement, which
we contributed to the Volcano scheduler \cite{volcano,volcanopr}, discovers
GPU-to-NUMA mappings and co-locates CPU and GPU on the same NUMA node. Combining
the two---utilization-driven scaling from \texttt{keda-gpu-scaler} and
topology-aware placement from Volcano---addresses both scaling and placement.

\section{Evaluation}

\subsection{Metric latency}
Table~\ref{tab:latency} compares the end-to-end metric latency and component count
of the pipeline approaches with the direct-NVML approach.

\begin{table}[h]
\centering
\begin{tabular}{L{0.42\textwidth} l l}
\toprule
\textbf{Approach} & \textbf{Metric latency} & \textbf{Components} \\
\midrule
DCGM $\rightarrow$ Prometheus $\rightarrow$ HPA & 15--30~s & 5 \\
DCGM $\rightarrow$ Prometheus $\rightarrow$ KEDA & 10--20~s & 4 \\
\texttt{keda-gpu-scaler} (direct NVML) & 2--4~s & 2 \\
\bottomrule
\end{tabular}
\caption{Metric latency and component count.}
\label{tab:latency}
\end{table}

\subsection{Production observations}
We have run \texttt{keda-gpu-scaler} on a multi-node A100 cluster serving LLM
inference with vLLM. We observed the following. First, device memory is the right
signal for vLLM: SM utilization stayed near 60--70\% even under heavy load
because vLLM batches requests, whereas VRAM pressure tracked request-queue depth
closely, with latency spiking within seconds once the KV cache filled past
roughly 80\%. Second, scale-to-zero produced meaningful savings: development and
staging endpoints sat idle for most of each day, and releasing those GPUs cut
non-production GPU spend substantially. Third, the short poll interval
mattered: with the Prometheus pipeline a traffic burst could saturate GPUs for
15--30~s before a new pod scheduled, whereas direct NVML surfaced the spike within
one poll cycle. Fourth, the best aggregation strategy is workload-dependent:
\texttt{avg} suited tensor-parallel serving across four GPUs, while \texttt{max}
caught a hot device faster when several smaller models shared a node. These are
observations from a single deployment rather than a controlled benchmark, and
results will vary with model size, request patterns, and node configuration.

\subsection{Operational overhead}
Compared with a DCGM-plus-Prometheus pipeline, the DaemonSet adds no cluster-wide
components, keeps all state in memory rather than in a time-series database, and
is configured through \texttt{ScaledObject} YAML rather than PromQL queries,
scrape configs, and adapter rules.

\section{Application: Agentic AI and Autonomous Networks}

Beyond conventional inference serving, GPU scale-to-zero is valuable wherever
demand is bursty and driven by discrete, high-severity events. Autonomous
(``self-healing'') telecom networks are one such setting \cite{comsoc}. Agentic
reasoning over a network failure requires loading multi-gigabyte foundation-model
weights into device memory, which is prohibitively expensive to keep running
continuously ``just in case'' an anomaly occurs. Two properties make standard HPA
a poor fit here: reactive CPU thresholds incur a cold-start penalty that can
breach sub-second service-level objectives, and hardware utilization is a
misleading signal because a GPU can be busy on a trivial task while a critical
event waits in the queue.

Coupling event-driven activation---for example, the depth or severity of an
anomaly queue drawn from the network data plane---with GPU scale-to-zero lets
operators provision cognitive compute only during active reasoning and release it
immediately afterward. This separates a lightweight, CPU-based observability
control plane (the ``watchers'') from a GPU-backed execution plane (the
``thinkers'') that defaults to zero instances, and reduces both cost and the
idle-power carbon footprint. In this architecture \texttt{keda-gpu-scaler}
supplies the GPU-aware activation signal that lets the orchestrator scale
cognitive workloads on hardware state rather than on coarse CPU metrics.

\section{Related Work}

DCGM Exporter exposes NVIDIA DCGM metrics to Prometheus but requires the full
pipeline and offers no direct KEDA integration. GPU feature discovery labels
nodes with static GPU properties but no runtime metrics. HAMi \cite{hami}
focuses on GPU sharing and partitioning rather than autoscaling. Higher-level AI
workload managers typically rely on standard HPA. Node autoscalers such as
Karpenter scale nodes rather than pods and are not GPU-metric aware.
\texttt{keda-gpu-scaler} is complementary: DCGM Exporter remains useful for
dashboards, while the scaler moves the scaling \emph{decision} off the monitoring
pipeline and onto a direct hardware read.

\section{Future Directions}

Planned directions include AMD ROCm support using a \texttt{rocm-smi}-based
collector behind the same interface, per-partition MIG metrics, NVLink
topology-aware scaling, reading vLLM queue depth directly from the engine for
predictive scaling, and cross-node (fleet-wide) metric aggregation for
cluster-level scaling decisions.

\section{Conclusion}

Kubernetes has no visibility into GPU state, and existing workarounds add latency
and operational complexity. A DaemonSet that reads NVML directly and serves KEDA
over gRPC respects the underlying constraints---NVML requires cgo, GPU devices
are node-local, and GPU tooling evolves faster than the KEDA release cycle---and
reduces metric latency to a few seconds while enabling GPU scale-to-zero. The
software is open source and available at \cite{kgs}.

\begin{thebibliography}{99}

\bibitem{keda} KEDA: Kubernetes Event-driven Autoscaling. \url{https://keda.sh}

\bibitem{kgs} P. Madduri. keda-gpu-scaler: GPU-aware autoscaling for Kubernetes
from real hardware metrics. Zenodo, 2026. \url{https://doi.org/10.5281/zenodo.22866676}

\bibitem{nvml} NVIDIA Corporation. NVIDIA Management Library (NVML).
\url{https://developer.nvidia.com/management-library-nvml}

\bibitem{gonvml} NVIDIA Corporation. go-nvml: NVML bindings for Go.
\url{https://github.com/NVIDIA/go-nvml}

\bibitem{keda-extscaler} KEDA External Scaler specification.
\url{https://keda.sh/docs/latest/concepts/external-scalers/}

\bibitem{keda7538} KEDA issue \#7538: GPU scaler discussion.
\url{https://github.com/kedacore/keda/issues/7538}

\bibitem{vllm} W. Kwon, Z. Li, S. Zhuang, Y. Sheng, L. Zheng, C. H. Yu,
J. E. Gonzalez, H. Zhang, and I. Stoica. Efficient Memory Management for Large
Language Model Serving with PagedAttention. In \emph{Proc. 29th ACM Symposium on
Operating Systems Principles (SOSP)}, 2023.

\bibitem{volcano} Volcano: Cloud-native batch scheduling system.
\url{https://volcano.sh}

\bibitem{volcanopr} P. Madduri. GPU NUMA-aware scheduling. Volcano PR \#5095.
\url{https://github.com/volcano-sh/volcano/pull/5095}

\bibitem{hami} HAMi: Heterogeneous AI Computing Virtualization Middleware.
\url{https://github.com/Project-HAMi/HAMi}

\bibitem{comsoc} P. Madduri and A. L. Thakur. The Financial Trap of Autonomous
Networks: Scaling Agentic AI in the Telecom Core. IEEE ComSoc Technology Blog,
March 2026.
\url{https://techblog.comsoc.org/2026/03/30/the-financial-trap-of-autonomous-networks-scaling-agentic-ai-in-the-telecom-core/}

\end{thebibliography}

\end{document}
